% ============================================================
%  REWRITTEN: Related Work
%  Target: JAIR
%  Changes from AAAI version:
%   - Expanded from 3 short paragraphs to 5 subsections (~3 pages)
%   - Adds coverage of WMC competition solvers (E4 context)
%   - Adds Goral et al. 2024, Onaka et al. 2025 (AI Reviewer request)
%   - Clearly distinguishes our approach from compositional WMC
%   - Adds ontology-based prediction subsection
% ============================================================

\section{Related Work}
\label{sec:related-work}


% ────────────────────────────────────────────────────────────
\subsection{Constrained and Ontology-Based Prediction}
% ────────────────────────────────────────────────────────────

Hierarchical multi-label classification (HMC) requires that positive
predictions respect a taxonomic parent-child ordering: a label assigned
to any node must also be assigned to all its ancestors
\citep{silla2011survey}. Methods for enforcing this constraint range
from architecture design---encoding hierarchy into the classifier
structure \citep{giunchiglia2020chmcnn}---to constraint-aware losses
such as semantic loss \citep{xu2018semanticloss} and SPL \citep{ahmed2022spl},
to posterior regularization \citep{ganchev2010posterior} and
constrained decoding \citep{diligenti2017semantic}.
Correction methods applied at inference time to fixed predictor outputs
include upward max-closure, which propagates positive predictions
upward to all ancestors \citep{obozinski2008consistent}, and
sum-product belief propagation over the constraint graph
\citep{giunchiglia2020chmcnn}.

In protein-function prediction, methods such as DeepGO
\citep{kulmanov2018deepgo}, DeepGOZero \citep{kulmanov2022deepgozero},
and DeepGOCNN \citep{kulmanov2020deepgoplus} assign per-term scores
evaluated under the CAFA framework \citep{radivojac2013cafa,zhou2019cafa3}.
Post-processing methods---including max-closure and loopy BP---are
applied to improve constraint satisfaction after training. Our work
studies WMC as a plug-in inference layer in this setting, targeting
the same fixed-score correction objective while returning a fully
normalized posterior rather than a constraint-satisfying point estimate.


% ────────────────────────────────────────────────────────────
\subsection{Exact Weighted Model Counting}
% ────────────────────────────────────────────────────────────

Weighted model counting computes the sum of weights of satisfying
assignments of a Boolean formula \citep{chavira2008wmc}. It generalizes
model counting (\#SAT) to real-valued weights and underlies probabilistic
inference in graphical models, Bayesian networks, and probabilistic
logic programs \citep{fierens2015problog,de2007probfoil}.

Knowledge compilation translates a formula into a tractable circuit
representation that supports WMC in linear time in the circuit size.
Representative target languages include binary decision diagrams (BDDs),
deterministic decomposable negation normal form (d-DNNF)
\citep{darwiche2002kcmap}, and sentential decision diagrams (SDDs)
\citep{darwiche2011sdd}. Compilation cost is exponential in treewidth,
and the resulting circuit may be exponential in the formula size in the
worst case. Variable elimination approaches such as ADDMC
\citep{dudek2020addmc}, DPMC \citep{dudek2020dpmc}, and
sharpSAT-TD process the formula directly, avoiding a separate
compilation step. The d4 compiler \citep{lagniez2017improved}
interleaves clause learning with compilation and performs competitively
on industrial instances. These solvers are evaluated annually at the
Model Counting Competition \citep{fichte2021mcc}; we include ADDMC,
DPMC, and d4 in our experimental comparison.

All of these methods treat the full variable set monolithically: they
do not exploit ontology-specific structure (sparse DAG, large scale,
predictor-valued weights) to partition computation. The regime targeted
by ontology-conditioned WMC---hundreds of thousands of variables,
moderately simple logical structure, real-valued weights changing at
each inference call---is distinct from the small-to-medium Boolean
instances that dominate benchmark suites.


% ────────────────────────────────────────────────────────────
\subsection{Compositional and Modular Model Counting}
% ────────────────────────────────────────────────────────────

Compositional WMC decomposes a formula into independent or loosely
coupled components and recombines their counts, reducing cost when
the formula is nearly disconnected \citep{dal2017partitioned,
dal2021compositional}. Dal et al.\ \citeyearpar{dal2021compositional}
generalize this to approximate probabilistic inference in Bayesian
networks, communicating between components through separator variables
with a product approximation. Ontology modularity, studied in the
knowledge-representation community, identifies logically meaningful
decompositions based on the import closure or atomic decomposition of
a description logic ontology \citep{cuencagrau2008modular,
delvescovo2011atomic}; these decompositions are defined at the
ontology level and do not account for predictor-weighted inference.

Our framework specializes compositional WMC to predictor-weighted
ontology theories. Three aspects distinguish it from prior compositional
approaches. \emph{First}, the module decomposition is driven by the
structural width of the constraint graph under the given predictor
weights, not by logical or semantic criteria---the core-selection
algorithm responds to the elimination order induced by the weighted
theory at inference time. \emph{Second}, our separator factors are
\emph{joint} over all boundary variables, retaining the full dependence
that the eliminated module induces on the rest of the ontology; prior
compositional approaches use factored or independence approximations
at wide separators. \emph{Third}, we provide a practical construction
for wide-boundary separators via oracle queries and TT fitting---a regime
($|B_j|$ in the hundreds to thousands) not addressed by existing
compositional methods. Dal et al.\ \citeyearpar{dal2021compositional}
explicitly note that their approach is limited to small separators.


% ────────────────────────────────────────────────────────────
\subsection{Tensor-Based Inference and Model Counting}
% ────────────────────────────────────────────────────────────

Low-rank tensor representations support efficient operations on
high-dimensional probability distributions. Tensor-train (TT)
decompositions \citep{oseledets2011tt} and cross-approximation
algorithms \citep{oseledets2010ttcross} provide compact representations
of functions over exponentially many inputs. In graphical models,
tensor-network contractions have been used for marginal inference
\citep{novikov2014mrf} and as an alternative to belief propagation
\citep{wrigley2017tensorbp}. Ducamp et al.\ \citeyearpar{ducamp2020lowrank}
apply low-rank decompositions to potential functions in Bayesian networks.

In model counting, Goral et al.\ \citeyearpar{goral2024semiring} introduce
semiring extensions that encompass WMC and weighted counting variants,
establishing algebraic foundations for combining tensor operations with
counting. Onaka et al.\ \citeyearpar{onaka2025ttkc} directly compare
tensor-train representations with ordered binary decision diagrams (OBDDs)
for knowledge compilation, finding that TTs achieve competitive size on
certain formula families. These works compile a static formula into a TT
or tensor-network representation. Our framework differs in two respects:
we do not compile the full formula but query the separator factor
pointwise using a boundary-clamped WMC oracle, and we project logical
support separately so that the TT fit is required only over the
interior of the feasible set---reducing the dynamic range of the fitted
function and making low-rank approximation more effective.

The use of nonnegative TT decompositions with alternating least-squares
fitting \citep{lee2001nmf} ensures that the separator factor represents
a nonnegative function, consistent with its interpretation as a
(unnormalized) marginal likelihood. This nonnegativity constraint, absent
in standard TT-SVD decompositions, is important for preserving the
probabilistic semantics of the WMC posterior.


% ────────────────────────────────────────────────────────────
\subsection{Message Passing and Approximate Inference}
% ────────────────────────────────────────────────────────────

Belief propagation (BP) \citep{pearl1988probabilistic} is exact on
trees and provides an efficient but approximate alternative to junction-tree
inference on graphs with cycles. On ontology graphs, loopy BP passes
messages along is-a edges, but reconvergent paths cause double-counting
of shared evidence; the resulting fixed-point marginals differ from
the globally normalized WMC posterior (Section~\ref{sec:local-reconvergence}).
Damped and synchronous BP variants \citep{murphy1999loopy} improve
convergence but do not correct the fundamental reconvergence error.
Our framework retains exact message passing where the structure permits
(inside each module and over the reduced core) while using joint
separator factors to handle reconvergent interfaces between modules.

Expectation propagation \citep{minka2001ep} and variational methods
\citep{wainwright2008graphical} provide alternative approximation
schemes but do not directly minimize the error in the WMC normalizer.
For fixed predictor scores and hard logical constraints, the
WMC posterior is the unique target, and any approximation is measured
by its deviation from that posterior---the criterion tracked by our
TV-distance stability bound (Proposition~\ref{prop:boundary-error}).


% ────────────────────────────────────────────────────────────
% ---- BIBLIOGRAPHY ENTRIES TO ADD (not in original .bib) ----
%
% @inproceedings{goral2024semiring,
%   author    = {Goral, Andreas and Giesen, Joachim and Blacher, Mark and
%                Staudt, Christian and Klaus, Jonas},
%   title     = {Model Counting and Sampling via Semiring Extensions},
%   booktitle = {Proceedings of the AAAI Conference on Artificial Intelligence},
%   volume    = {38},
%   number    = {18},
%   pages     = {20395--20403},
%   year      = {2024}
% }
%
% @inproceedings{onaka2025ttkc,
%   author    = {Onaka, Ryoma and Nakamura, Kengo and Nishino, Masaaki and
%                Yasuda, Norihito},
%   title     = {Tensor Decomposition Meets Knowledge Compilation:
%                A Study Comparing Tensor Trains with {OBDDs}},
%   booktitle = {Proceedings of the AAAI Conference on Artificial Intelligence},
%   volume    = {39},
%   number    = {14},
%   pages     = {15109--15117},
%   year      = {2025}
% }
%
% @article{lagniez2017improved,
%   author    = {Lagniez, Jean-Marie and Marquis, Pierre},
%   title     = {An Improved Decision-{DNNF} Compiler},
%   journal   = {Proceedings of the International Joint Conference on
%                Artificial Intelligence},
%   pages     = {667--673},
%   year      = {2017}
% }
%
% @inproceedings{fichte2021mcc,
%   author    = {Fichte, Johannes Klaus and Hecher, Markus and Hamiti, Florim},
%   title     = {The Model Counting Competition 2020},
%   booktitle = {ACM Journal of Experimental Algorithmics},
%   volume    = {26},
%   year      = {2021}
% }
%
% @article{lee2001nmf,
%   author    = {Lee, Daniel D. and Seung, H. Sebastian},
%   title     = {Algorithms for Non-negative Matrix Factorization},
%   journal   = {Advances in Neural Information Processing Systems},
%   volume    = {13},
%   year      = {2001}
% }
%
% @inproceedings{dal2017partitioned,
%   [already in references.bib as dal2017]
% }
% ---- END OF BIB ENTRIES ----
